\section{Troubleshooting}
\label{sec:Troubleshooting}

Most operational failures are caused by startup timing, missing topics,
incomplete task logic, Kafka connectivity, or confusion between execution state
and monitoring state. Start with the runtime signals before assuming a generator
defect.

\subsection{Kafka UI logs connection failures at startup}

This is usually normal in the local setup. Kafka UI often starts before the broker is fully
ready and logs transient connection failures. If the broker then comes up and the errors stop,
this is not itself evidence of a Durga runtime failure.

\subsection{Monitoring app says the state store is not queryable yet}

This normally means Kafka Streams has started but the state stores are not yet ready for queries.
The correct response is usually to wait briefly and retry rather than to assume the topology is
broken immediately.

In production, repeated or long-lived readiness failures should be treated as an
operational incident. Check broker reachability, lifecycle-topic consumer lag,
Streams application id stability, changelog topic health and whether a deployment
started with a fresh state directory or missing permissions.

\subsection{A generated process compiles but does not advance}

The most common causes are:
\begin{itemize}
\item topics were not created,
\item the relevant generated application is not running,
\item an XOR branch condition still contains placeholder logic,
\item a user or manual task entered a waiting state and needs explicit completion,
\item the process instance was cancelled by a boundary or interrupting event subprocess,
\item the worker processed a record but failed before committing its offset, causing
      a replay that requires idempotent handler logic.
\end{itemize}

\subsection{Boundary behavior does not appear to trigger}

Check these points first:
\begin{itemize}
\item the triggering lifecycle event is actually emitted,
\item the boundary is interrupting or non-interrupting as intended,
\item the attached task or subprocess reached the expected state,
\item failure or escalation was emitted through the generated helper path rather than by bypassing
the lifecycle model.
\end{itemize}

\subsection{An external completion arrives after cancellation}

The generated runtime now suppresses much of this post-cancel progress, but late events can still
be confusing during debugging. In such cases, the canonical \texttt{process-events-\{processId\}} stream is the
best source of truth: check whether cancellation happened before the late completion arrived.

\subsection{Monitoring counts do not match raw topic intuition}

This is usually a model misunderstanding rather than a bug. Execution topics contain history and
routing traffic. The monitoring counts are derived from the latest projected instance state, which
is the intended read model for ``how many instances are in state X right now''.

\subsection{Dashboard trend or latency data looks incomplete}

The dashboard reads materialized monitoring stores, not raw execution topics. Missing trend or
latency rows usually means no supported lifecycle event has reached the monitoring topology yet,
the Streams stores are still restoring, or old monitoring topics contain records produced by an
older projection version. For persistent inconsistencies, compare \texttt{process-events-\{processId\}},
\texttt{process-trends}, \texttt{process-latency}, and the monitoring application logs.

\subsection{The generated output looks harder to understand than expected}

This often comes from BPMN naming. Unclear BPMN ids produce unclear Java classes, topics and helper
artifacts. Cleaning up the BPMN model identifiers usually improves the generated output
more than post-processing the generated files.
